\begin{afterword}
\summaryhead

The post introduces the basic objects of the sequence. It describes reality as a program. The program holds configurations, such as ``a photon heading toward A,'' and each configuration stores a complex number called an amplitude. A half-silvered mirror passes the amplitude on multiplied by 1 along the straight path and by $i$ along the turned path. Amplitudes cannot be seen directly: counting detector clicks gives only the ratios of their squared moduli.

With one half-mirror, both detectors click equally often. In an interferometer with two half-mirrors and two full mirrors, the amplitudes arriving at one detector cancel, and every photon reaches the other. Blocking either path makes both detectors click equally again. A photon that simply went one way or the other could not do this, and neither could one that the mirrors marked to turn differently next time.

Probabilities are never negative, so they cannot cancel; therefore amplitudes are not degrees of belief. Configurations are not possible worlds, and what could have happened but did not has no effects. Amplitudes have effects, so they are real. Real configurations cover every particle in the universe and form a continuous space.

\responsehead

The post's physics is right where it counts. Every worked number checks: the amplitudes at each stage, the cancellation at E, the even split when a path is blocked, and the probability example. Its central claim, that amplitudes arriving by two paths can cancel and that no particle going one way or the other under fixed rules can reproduce this, is the standard lesson of the Mach--Zehnder interferometer.

Its simplifications are mostly flagged, but two affect what a reader takes away. The mirror rule leaves out the factor $1/\sqrt{2}$ of the standard beam splitter, so the squared moduli double at each half-mirror. The post's reported results survive, because it compares only configurations that passed the same number of mirrors. But with a path blocked, half of all photons stop at the block, so ``Detector 1 goes off half the time'' holds only for the photons that reach a detector. And the experiments shown do not single out a ratio of $i$ to 1; the editor's note at the end gives another rule that fits them.

The larger step is from physics to interpretation, which the post takes in one move. Amplitudes have visible effects, it says, so they are causes, and so they are real, ``out there in the territory.'' Against one rival, the view that the quantum state is knowledge about some underlying physical state, later work supports the post: a 2012 theorem by Pusey, Barrett and Rudolph rules that view out under an independence assumption. Other rivals are untouched. The argument that amplitudes are not degrees of belief does not reach QBism, which says that the quantum state as a whole specifies an agent's credences about outcomes, not that each amplitude is one. The argument against billiard balls does not reach pilot-wave theory, in which a particle takes one path, guided by a wave on both, with the same predictions. What the experiments show is that amplitudes are needed, which no interpretation denies. That they are real causes is the disputed premise, and the post asserts it.

Two remarks about the final picture go beyond what the post establishes. The rule that click counts follow squared moduli is postponed; later in the book the author writes that ``There is no known reason for the Born probabilities.'' And configurations over ``joint positions of all the particles in the universe'' belong to quantum mechanics with a fixed number of particles; the post's photons are described by quantum field theory.

In short: a correct introduction to amplitudes and interference, whose worked examples all check, which treats the realist reading of amplitudes as something the experiments show, when they show only that amplitudes are needed.
\end{afterword}
